\chapter{Background} \label{chap:background}
\section{What are network sockets} \todo{Headers should be removed after editing is finished}
Anyone perfoming any kind of low-level network programming in a UNIX
environment are bound to have encountered sockets. Sockets exist
in several forms, and can be used for both inter-process communication
and network communication. 
In its bost basic definition, a network sockets is an abstraction
for a network connection, through which an application can send
and receive data \cite[p. 8]{donahooTCPIPSockets2009}. 
First mentioned in 1970 in relation to ARPANET, the
simplicity of sockets and their ubiquity in all UNIX distributions
since version 4.2 of the Berkeley Software Distribution have
led to them becoming the de facto standard for low-level network
programming \cite{crockerNewHostHostProtocol1970}\cite[p. 502]{tanenbaumComputerNetworks2011}.

Modern sockets are standardized across UNIX-like systems through
the POSIX standard. Implementations of the C POSIX library define
the socket API primarily through the \texttt{sys/socket.h} header file.

When creating a socket on an UNIX-like system, the application
invokes the \texttt{socket} system call. This system call
\say{creates an endpoint for communication and returns a file 
descriptor that refers to that endpoint} \cite{Socket2LinuxManual}.
This system call takes three parameters, the domain, the type
and the protocol which will be used for the socket.
The domain specifies the type of
addressing which will be used, for example \texttt{AF\_UNIX}
for local sockets, \texttt{AF\_INET} for IPv4 and \texttt{AF\_INET6}
for IPv6. The type specifies the protocol which will be used over the socket, for
example \texttt{SOCK\_DGRAM} for datagrams and \texttt{SOCK\_STREAM}
for a reliable stream of bytes.
The protocol parameter is used if there are multiple protocols
which support the requirements given by the domain and type.

The fact that the application has to specify both the domain
and the type of the socket means that both the
underlying addressing and the protocol are bound at compile time.
This means that barring fairly advanced logic when setting up
the necessary sockets, the application will be bound to a specific
protocol and addressing scheme. Therefore, updating
the underlying protocol will require a recompilation of the application. 
Naturally, this makes adopting new protocols difficult, even
if the new protocol is better suited to the application's requirements.

Another challenge with new transport-layer protocols apart
from TCP and UDP is that developing them has proven to be challenging, 
in part because of TCP and UDP's ubiquity. 

Network middleboxes have have been
optimized to support TCP and UDP, and may sometimes
even drop packets if they contain another transport layer protocol.
The resulting difficulty in innovating on the transport-layer
due to network growing accustomed to the current state of things
is known as network ossification. Due to this network ossification,
most new protocols are built on top of UDP.

Due to UDP's simplicity, it presents very little overhead
and applications can themselves implement needed features
such as reliability and congestion control.
This has become so common that
it has prompted an official RFC from the IETF, describing
the best practices for implementing congestion control on top of UDP 
\cite{eggertUnicastUDPUsage2008}.

\section{What is TAPS}
Solving the problem with the strong coupling between application
level code and the underlying transport protocol
is one of the key goals of TAPS. This will allow faster adaptation
of new protocols and abstract away network-layer details from the
application layer. From the application's point of view, this
means that it will be able to gain access to new protocols, such
as QUIC, without having to implement its own fallback logic
in cases where QUIC is not available. It will also simplify
the networking logic, tying it more closely to how modern
applications actually use the network.

Transport Services (TAPS) is a \say{application Programming Interface 
(API) to the transport layer that enables the selection of 
transport protocols and network paths dynamically at runtime}
\cite{trammellAbstractApplicationProgramming2025}. TAPS is
designed to remove the compile-time binding to transport layer
protocols, enabling updates to underlying transport protocols
without having to rewrite application logic. TAPS does not
require that both endpoints of a connection use TAPS, and TAPS
can communicate with any endpoint which supports the underlying
transport protocol being used.

The following figure explains the API model when using
sockets: \todo{Should these be here? If yes, recreate in tikz}

\begin{figure}[H]
    \centering
    \begin{lstlisting}
        +-----------------------------------------------------+
        |                    Application                      |
        +-----------------------------------------------------+
                |                 |                  |
          +------------+     +------------+    +--------------+
          |  DNS Stub  |     | Stream API |    | Datagram API |
          |  Resolver  |     +------------+    +--------------+
          +------------+          |                  |
                            +---------------------------------+
                            |    TCP                UDP       |
                            |    Kernel Networking Stack      |
                            +---------------------------------+
                                            |
        +-----------------------------------------------------+
        |               Network-Layer Interface               |
        +-----------------------------------------------------+
    \end{lstlisting}
    \caption{API model for socket programming \cite{trammellAbstractApplicationProgramming2025}}
\end{figure}

The following figure explains the API model when using TAPS:

\begin{figure}[H]
    \centering
    \begin{lstlisting}
        +-----------------------------------------------------+
        |                    Application                      |
        +-----------------------------------------------------+
                                  |
        +-----------------------------------------------------+
        |              Transport Services API                 |
        +-----------------------------------------------------+
                                  |
        +-----------------------------------------------------+
        |          Transport Services Implementation          |
        |  (Using DNS, UDP, TCP, SCTP, DCCP, TLS, QUIC, etc.) |
        +-----------------------------------------------------+
                                  |
        +-----------------------------------------------------+
        |               Network-Layer Interface               |
        +-----------------------------------------------------+
    \end{lstlisting}
    \caption{API model for TAPS \cite{trammellAbstractApplicationProgramming2025}}
\end{figure}

\subsection{TAPS abstractions}
Abstracting away the underlying transport protocols is done by
presenting a set of abstractions to the application. The most
central of these abstractions is the Connection. A Connection
represents a \say{shared state of two or more endpoints that persists
across messages that are transmitted between these endpoints}
\cite{welzlUsageTransportFeatures2018}.
When working with Connection objects, TAPS provides five high
level \textit{groups of actions} that an application can perform:
\begin{itemize}
    \item \texttt{Preestablishment}: Pass properties to specify
        requirements and preferences for establishing a Connection.
    \item \texttt{Establishment}: Prepare a Connection object for data
        transfer.
    \item \texttt{Data transfer}: Sending and receiving messages
        over the Connection.
    \item \texttt{Event handling}: Notifications that an application
        can receive from a Connection, and the handling of them.
    \item \texttt{Termination}: How transmission is stopped and
        a Connection is torn down.
\end{itemize}

In addition, the lifeline of a Connection is divided into three distinct
phases, \textit{preestablishment}, the \textit{Established state}, and 
\textit{termination of a Connection} \cite{paulyArchitectureRequirementsTransport2025}.
Each of these phases are associated
with one or more of the \textit{groups of actions} that an application can perform,
as described above. 

After the desired transport properties and cryptographic parameters have
been set during the preestablishment phase, a connection can be established. 
When a Connection is in the established phase the Connection can be used to
transfer data. 

There are three ways of
establishing a connection: \textit{initiate}, \textit{listen} and
\textit{rendezvous} \cite{paulyArchitectureRequirementsTransport2025}.

\textit{Initiate} is used to establish a connection to a remote endpoint which
is assumed to be listening for incoming connections. How this connection
is established under the hood, is dependent on the underlying transport
protocol, which has been selected based on the specified transport properties.

\textit{Listen} is used to passively open a connection. This is done in order to
wait for connections from remote endpoints. Listen returns a listener
object. When a remote endpoint connects to this listener, a \texttt{ConnectionReceived}
event occurs. This event contains a Connection object for the 
established connection.

\textit{Rendezvous} is used to establish a peer-to-peer connection between two endpoints
simultaneously.

After a Connection has been established, it can either
be closed or aborted, with closing being associated with normal termination,
and aborting being abnormal termination associated with errors.

If a single application has multiple connections, these can be organized
into connection groups. Connection groups represent a set of connections
which shared connection properties and state generated by the underlying
transport protocols. Connections in the same connection group can be
multiplexed together, in order to more efficiently utilize the available
bandwith.

In addition to the Connection, TAPS also provides several other objects
abstracting key elements. For example Endpoint objects and Message objects.

\subsection{Selecting the transport protocol}
TAPS selects the underlying transport protocol by
letting the application specify a set of requirements
and preferences for the Connection behavior. This set of requirements
and preferences are known as \textit{Transport Properties}. These transport
properties can again be divided into three categories: \textit{selection
properties}, \textit{connection properties} and \textit{message properties}. 

\begin{itemize}
    \item \textit{Selection
properties} are properties used to select the path and protocol stack
to use for a given connection. 

    \item \textit{Connection properties} are used
during Connection establishment and to define the behavior of the
established Connection. 

    \item \textit{Message properties} are settings set on an
individual message. 
\end{itemize}

Typical transport properties are specified as
a \textit{preference}, with the possible values being \textit{require},
\textit{prefer}, \textit{no preference}, \textit{avoid} and \textit{prohibit}.
\textit{Require} means that the Connection establishment will fail if the property cannot
be satisfied, while \textit{prefer} will allow the application to continue.
This distinction is mirrored in the difference between \textit{prohibit} and \textit{avoid}.

An example of a selection property is Reliable Data Transfer.
This property is used to specify whether or not the application
should use a protocol which ensures that all data is received at 
a remote endpoint, in order, without loss or duplication. Setting
this to \textit{require} will force TAPS to use an underlying protocol
which guarantees this, such as TCP or QUIC. Setting this as
\textit{prohibit} will force TAPS to use a protocol which does not
guarantee this, such as UDP.


An example of a connection property is Timeout for Aborting Connection.
This property is used to specify the how long to wait before deciding that an
active connection has failed when trying to reliably deliver data to the Remote Endpoint.
This naturally can only take effect if the underlying protocol supports reliability.

An example of a message property is \texttt{priority}. This property specifies
the prority of a given message relative to other messages over the same connection.
Lower priority message will yield in favor of higher priority messages.

TAPS is designed to provide an interface which more accurately
reflects how a typical application uses the network. Socket programming
is for example blocking by default, while most modern applications
tinteract with the network in an asynchronous way 
\cite{trammellAbstractApplicationProgramming2025}.

An additional
example of how TAPS is designed to reflect how applications
use the network is that TAPS is message-oriented. 
Being stream-oriented or message-oriented has to do with how the
transferred data is represented. In a stream oriented protocol
such as TCP, the data is represented as a stream of bytes. From
the application's point of view, the data arrives continuously
and not in discrete chunks. This means that the application
cannot easily differentiate between different messages. 

TCP has to
handle  data arriving in discrete IP packets and glue them together into
a stream of bytes. This also includes handling messages arriving
out of order, which has to be handled before they can be
presented as a stream in order.
Since most applications
use the network in a message-oriented way, they
have to include delimiters and other logic to separate TCP 
data into discrete messages. For example, several protocols
such as HTTP contain length fields, which are used to delimit
transferred data. 

This is the reason why TAPS is designed to be message-oriented,
as it simplifies the logic needed for most applications to
use the network in the desired way. Despite being message-oriented,
TAPS can still support applications reliant on 
streaming through the use of \textit{framers}.

A framer is a data translation layer, which is used to define
how encapsulate or encode outbound message, and how to decapsulate
or decode inbound messages. This for example allows applications
to receive a byte stream as input if needed.

TAPS is backwards compatible with existing protocols as it does not
define how these underlying protocols should transfer data over
the network. TAPS is also designed to coexist with existing
socket programming APIs, allowing a gradual transition from
socket programming to TAPS.

Another key feature difference between TAPS and sockets is that socket
APIs are commonly limited to communicating with a single address over 
a single interface. TAPS is deisgned to handle multiple endpoints,
protocols and paths.


\section{What is quic}

QUIC is a internet protocol developed by Google, released to the
public in 2013 described in RFC 9000
\cite{iyengarQUICUDPBasedMultiplexed2021}. QUIC is a transport
layer protocol which is designed to be a replacement for TCP.
QUIC is a connection-oriented protocol which is built on top of
UDP. QUIC enables 0-RTT connection establishment and integrates
encryption as a mandatory feature of the protocol.

The abstraction presented to an application for communication
via QUIC is called a stream. A stream is an either either
unidirectional or biderectional ordered stream of data, much
like what an application sees when using TCP. Streams can be created and
closed by either endpoint, and multple streams can exist over a
single QUIC connection. This is an example of the concept of
multiplexed connections. Individual pieces of data sent over a
stream are encapsulated as stream frames. The stream frame boundaries
are not expected to be preserved, and are therefore not visible
to the application.

By using the stream ID and the byte offset of the stream frame,
the receiving endpoint reassembles the data into the correct
order. All implementations must support in-order delivery of
data, they may choose to allow out-of-order delivery of data as
well, but this is not a requirement of the protocol This only
happens on a per-stream basis, so when multiplexing multiple
streams over a single connection, the streams are independent of
each other and we avoid head-of-line blocking.

A given stream frame has the following format, as describted in
RFC 9000:

\begin{figure}[H]
	\centering
	\begin{lstlisting}
        STREAM Frame {
          Type (i) = 0x08..0x0f,
          Stream ID (i),
          [Offset (i)],
          [Length (i)],
          Stream Data (..),
        }
    \end{lstlisting}
	\caption{Format of a QUIC stream frame \cite{iyengarQUICUDPBasedMultiplexed2021}}
	\label{fig:stream_frame_format}
\end{figure}

A given stream can exist in several different states.
These states differ whether the stream is sending or receiving
data. For a stream sending data, the states are:
\texttt{Ready}, \texttt{Send}, \texttt{Data Sent}, \texttt{Reset Sent},
\texttt{Data Recvd}, and \texttt{Reset Recvd}.
A receiving streams states maps the perceived states of the
sending stream. This means that there are some states
which are not represented in the receiving stream, as they
cannot be inferred from the data received. The states of a
receiving stream are: 
\texttt{Recv}, \texttt{Size Known}, \texttt{Data Recvd}, \texttt{Reset Recvd},
\texttt{Data Read}, and \texttt{Reset Read}.
A bidirectional stream consists of both sending and receiving parts,
and can therefore exist in states belonging to both sending and
receiving streams as described previously.

\begin{figure}[H]
    \centering
    \begin{lstlisting}
       o
       | Create Stream (Sending)
       | Peer Creates Bidirectional Stream
       v
   +-------+
   | Ready | Send RESET_STREAM
   |       |-----------------------.
   +-------+                       |
       |                           |
       | Send STREAM /             |
       |      STREAM_DATA_BLOCKED  |
       v                           |
   +-------+                       |
   | Send  | Send RESET_STREAM     |
   |       |---------------------->|
   +-------+                       |
       |                           |
       | Send STREAM + FIN         |
       v                           v
   +-------+                   +-------+
   | Data  | Send RESET_STREAM | Reset |
   | Sent  |------------------>| Sent  |
   +-------+                   +-------+
       |                           |
       | Recv All ACKs             | Recv ACK
       v                           v
   +-------+                   +-------+
   | Data  |                   | Reset |
   | Recvd |                   | Recvd |
   +-------+                   +-------+
    \end{lstlisting}
    \caption{States of a QUIC stream sending data \cite{iyengarQUICUDPBasedMultiplexed2021}}
    \label{fig:stream_states}
\end{figure}

Another important abstraction presented by QUIC is the connection.
A QUIC connection is a shared state between a client and a server 
\cite{iyengarQUICUDPBasedMultiplexed2021}.
A connection is established through a handshake, with support
for 0-RTT handshakes as mentioned previously. The connection
is identified by a connection ID, which is selected by each
endpoint's peer. Packets sent to this peer are identified as
belonging to a given connection, by a destination connection ID
field. Connection IDs make addressing in QUIC independent of addressing in lower layer
protocols \cite{iyengarQUICUDPBasedMultiplexed2021}. Each endpoint
has a unique connection ID, making sure that external observers cannot
identify that the two endpoints are communicating with each
other by looking at the connection ID.

Using a connection ID allows the connectoin to survive a change
in the addressing of the underlying protocol \cite{iyengarQUICUDPBasedMultiplexed2021}.
This functionality is enabled by default, but can be disbled by
one of the peers. When an address is changed, each endpoint has to
independently verify that they can communicate using the new address.
Path validation happens using a special type of stream frame
known as \texttt{PATH\_CHALLENGE} and \texttt{PATH\_RESPONSE}
frames. The \texttt{PATH\_CHALLENGE} frame contains unpredictable
data which must be mirrored back in the \texttt{PATH\_RESPONSE} frame.
If the endpoint validation fails, which happens if no matching response
is received within a given time period, the endpoints can continue
communicating using other known addresses. If there are no valid
addresses, the affected endpoint can signal this to the other
peer by using a \texttt{NO\_VIABLE\_PATH} error.

Much like TCP, QUIC also implements congestion control, there are
however some key differences between the algorithms. A key concept
to QUIC's congestion control is the 'packet number space'. Packet
number spaces are \say{the context in which a packet can be processed
and acknowledged} \cite{iyengarQUICUDPBasedMultiplexed2021}. 
This context reflects the level of encryption and there 
are therefore three separate packet number spaces. Initial space,
for the initial handshake packets, handshake space, for packets
exchanged during the cryptographic handshake, and application
data space, for packets exchanged after the cryptographic handshake
has been completed. Separating these spaces allows QUIC to
avoid spurious retransmission from one space affecting another,
and allows QUIC to handle retransmissions more efficiently, for
example by using a different expected acknowledgement delay for packets sent
in the initial space, as these are not expected to be cached
by the peer. Although the packet number spaces are handled
separately when it comes to acknowledgements, they are still
unified when it comes to congestion control and round trip
time measurements.

Another difference between how QUIC and TCP handle acknowledgements
is that QUIC separates the transmission order from the delivery order.
This is done by using stream offset as delivery order, and packet number as
transmission order. This removes ambiguity about whether the
original or the retransmitted packet is acknowledged in the case
of retransmissions, making RTT measurements in these cases more accurate
\cite{iyengarQUICUDPBasedMultiplexed2021}. Another key difference
is that instead of TCP algorithms such as retransmission timeout (RTO)
and tail-loss probe (TLP), QUIC uses probe timeout (PTO).
PTO is an algorithm which is used to detect lost packets. Every
time a packet requiring an acknowledgement is sent, a timer is started. If
the timer expires, QUIC sends a probe packet and the PTO period is
set to twice its previous value. When an acknowledgement is received,
the PTO period is reset to its previous value, as long as the
connection has moved out of the Initial space. This is beneficial
as it can mitigate lost acknowledgements, and does not rely on
duplicate acknowledgements, making it beneficial in several cases
such as tail-loss. \todo{Make sure this is correct}

For QUIC a single connection can exist over several paths and
congestion control is done on a per-path basis. Multiple streams
can be transmitted over a single connection and data from
multiple stream frames can be transmitted in a single packet.
If a packet is lost, only applications using the streams contained
in that packet will be affected. This is a key difference from TCP,
where a single lost packet will block all other streams of data
over the same connection. This is called head-of-line blocking.
This problem is avoided by QUIC, since QUIC only guarantees
in-order delivery of data on a per-stream basis. This leads
to a better utilization of the available bandwidth, when
compared to multiplexing streams over a single TCP connection.

Future network ossification related to QUIC is also avoided
since QUIC mandates TLS 1.3 as part of the protocol.
By encrypting as much of the packet as possible, QUIC is able to
obscure details about the protocol from middleboxes, and
they cannot grow to depend on the contents of the packet.

As mentioned previously, one of many reasons behind the success
of TCP is
the universal support of sockets, which are implemented in kernel space.
There is no equivalent API for QUIC, and with a relatively
advanced implementation, users have to settle for one of the
existing libraries.


Currently, few applications can be built exclusively on top of
QUIC, as it requires support from both the client and the
server. This means that a client wanting to use QUIC, also
has to have support for fallback code, in case the server
it wants to connect to, does not have the required support.
As mentioned previously, making the adoption of new protocols
such as QUIC is one of the key selling points of TAPS. With TAPS,
the client would be able to specify the requirements for the
connection, and TAPS could select QUIC if it was supported, otherwise
falling back to TCP. This would require no change in logic for
the client, and would not require TAPS to be supported both
by the client and the server.

% =============== REVISED TO HERE ================

\section{TAPS configurations explained}

TAPS exposes several configurations which can be used to specify
a preference of requirement, which then sets a preference or a
requirement for which underlying transport protocol to use.
These configurations come in the form of transport properties.
These properties are broken down into three categories:
Selection properites, Connection properties, and message
properties. Selection properties are requirements for the
underlying transport protocol, which TAPS uses to select which
protocol to use for the connection. Connection properties are
properties which are set on a per-connection basis, for example
the timeout for aborting a connection or the connection
priority. Message properties are properties which are set on a
per-message basis, for example the priority or the lifetime of
the given message. A full list of properites can be found in the
Architecture and Requirements for Transport Services RFC
\cite{paulyArchitectureRequirementsTransport2025}.

Selection properties can, as mentioned above, be specified to
varying degrees of strictness. With five different values
\texttt{Require}, \texttt{Prefer}, \texttt{No Preference},
\texttt{Avoid}, and \texttt{Prohibit}. Here, \texttt{Require}
means that establishing the connection will fail if the property
cannot be satisfied with the available transport protocols.
\texttt{Prefer} means that the connection can fall back to a
protocol not supporting the property, if no protocol satisfying
the property is available. \texttt{Avoid} and \texttt{Prohibit}
function in the same way as \texttt{Require} and
\texttt{Prefer}, but selecting for protocols which do not have
the selected property. A simple example of a selection property
is selecting for reliable data transfer. If this was set as
required, TAPS would be forced to use for example TCP or QUIC as
the underlying protocol. If UDP was the only available protocol,
establishing the connection would fail. If the user selected
\texttt{prefer}, then UDP could be selected if nothing else was
available, even though it does not fulfill the property.

There are also more advanced selection properties available. Out
of these, the least intuitive are \texttt{Multistream
	Connections in Group}, \texttt{Full Checksum Coverage on
	Sending} and \texttt{Full Checksum Coverage on Receiving},
\texttt{Provisioning Domain Instance or Type}, \texttt{Use
	Temporary Local Address}, and \texttt{Notification of ICMP soft
	error message arrival}.

The \texttt{Multistream Connections in Group} selection property
allows the user \todo{TODO - select term for "user" of TAPS api
	(developer of networking application)} to select whether the
application would prefer multiple Connections within a
Connection Group to be provided by streams of a single
underlying transport connection where possible
\cite{paulyArchitectureRequirementsTransport2025}.
Multistreaming in this context is the act of providing multiple
streams of data over a single underlying connection. This is
advantageous as it allows for more efficient use of the network,
since the connection will better saturate the available
bandwidth. This avoids long convergence times for the connection
and allows fresh connection to utilize the large congestion
window already achieved, without having to wait for slow start
to ramp up. This is especially useful for applications with many
small streams of data, such as web browsers.

The \texttt{Full Checksum Coverage on Sending} and \texttt{Full
	Checksum Coverage on Re\-ceiving} selection properties specify
the application's need for protection against corruption for all
data transmitted or received on a given connection. This is used
to specify how many bytes of the message which needs to be
protected by a checksum. This checksum is used to protect the
data from corruption caused by lower-layer protocols. The
setting \texttt{Full coverage} means that the entire message is
protected, while \texttt{0} specifies no checksum is required.
If any other value than \texttt{Full coverage} is specified, the
connection may still fall back to \texttt{Full coverage}.

The \texttt{Provisioning Domain Instance or Type} selection
property allows the applicatio \todo{common application term}n
to control path selection by selecting for specific Provisioning
Domain or categories of domains
\cite{paulyArchitectureRequirementsTransport2025}. Provisioning
Domains define consistent sets of network properties that might
be more specific than network interfaces
\cite{paulyArchitectureRequirementsTransport2025}.

The \texttt{Use Temporary Local Address} selection property
allows the application \todo{common application term} to specify
whether or not to use temporary addresses which are a feature in
IPv6, used for greater privacy
\cite{gontTemporaryAddressExtensions2021},\cite{paulyArchitectureRequirementsTransport2025}.

The \texttt{Notification of ICMP soft error message arrival}
selection property lets an application \todo{common application
	term} specify whether or not it should be notified when an ICMP
error, which does not force termination of the connection,
arrives \cite{paulyArchitectureRequirementsTransport2025}.

In addition to specifying which underlying protocol should be
used through selection properties, the application \todo{common
	application term} using TAPS can manage the connections
themselves through transport properties. These can be specific
for a given transport protocol such as the selection properties,
or generic and be available for all transport protocols.
\todo[]{Should this be separated from all the above parts? are
	they all part of the same 'category'}

The \texttt{Connection Group Transmission Scheduler} transport
property decides which scheduler is used within a given
connection group, to divide the available capacity between the
connections. \cite{paulyArchitectureRequirementsTransport2025}.

The \texttt{Capacity Profile} transport property is used to
select a capacity profile, which will be a tradeoff between
delay, delay variation, and efficient use of the available
capacity based on the capacity profile specified
\cite{paulyArchitectureRequirementsTransport2025}.

The \texttt{Isolate Session} transport property set in order to
have new connections use as little cacheed information as
possible from previous connections not in the same connection
group \cite{paulyArchitectureRequirementsTransport2025}.
